\subsection{Orthogonal Isometries and Symmetry Groups}

Throughout, orthogonal geometries are finite-dimensional over a fixed field $\mathbb{F}$; characteristic restrictions are local.

\begin{proposition} \label{prop:orthogonal-isometry-preserves-quadratic-form} \hlt{Isometries from Quadratic Values}\\
Let $\mathcal{V}$ be orthogonal with $\operatorname{char}(\mathbb{F})\neq2$, put $q(\mathbf{x})\defeq\innerprod{\mathbf{x},\mathbf{x}}$, and let $T\in\operatorname{GL}(\mathcal{V})$. Then $T\in\mathcal{O}(\mathcal{V})$ if and only if $(\forall\,\mathbf{x}\in\mathcal{V})\,q(T\mathbf{x})=q(\mathbf{x})$.
\begin{proof}
\begin{enumerate}[label=\roman*.]
\setlength{\itemsep}{0pt}
\item[$(\Rightarrow)$] If $T$ is an isometry, then $(\forall\,\mathbf{x}\in\mathcal{V})\,q(T\mathbf{x})=\innerprod{T\mathbf{x},T\mathbf{x}}=\innerprod{\mathbf{x},\mathbf{x}}=q(\mathbf{x})$.
\item[$(\Leftarrow)$] Suppose that $q\circ T=q$. By linearity of $T$ and polarization, $(\forall\,\mathbf{x},\mathbf{y}\in\mathcal{V})$\\
$2\innerprod{T\mathbf{x},T\mathbf{y}}=q(T(\mathbf{x}+\mathbf{y}))-q(T\mathbf{x})-q(T\mathbf{y})=q(\mathbf{x}+\mathbf{y})-q(\mathbf{x})-q(\mathbf{y})=2\innerprod{\mathbf{x},\mathbf{y}}$. Cancellation of $2$ proves the form.
\end{enumerate}
\hlt{Structural Commentary}: Evaluate the form on the diagonal, then recover cross-pairings by polarization.\\
\hlt{Key Insight}: When $2$ is invertible, preserving every quadratic value is equivalent to preserving the bilinear form.
\end{proof}
\end{proposition}

\begin{definition} \label{def:orthogonal-symmetry-rotation-reflection} \hlt{Special Orthogonal Group and Symmetries}\\
Let $\mathcal{V}$ be nonsingular orthogonal with $\operatorname{char}(\mathbb{F})\neq2$, and put $q(\mathbf{x})\defeq\innerprod{\mathbf{x},\mathbf{x}}$.
\begin{enumerate}[label=(\roman*)]
\setlength{\itemsep}{0pt}
\item The \hlt{special orthogonal group} is $\operatorname{SO}(\mathcal{V})\defeq\{\sigma\in\mathcal{O}(\mathcal{V}):\det\sigma=1\}$; its elements are \hlt{rotations}.
\item For $q(\mathbf{u})\neq0$, the \hlt{symmetry} with normal line $\mathbb{F}\mathbf{u}$ is the map $\sigma_{\mathbf{u}}(\mathbf{x})\defeq\mathbf{x}-2\frac{\innerprod{\mathbf{x},\mathbf{u}}}{q(\mathbf{u})}\mathbf{u}$.
\end{enumerate}
\end{definition}

\begin{proposition} \label{prop:symmetry-matches-equal-length-nonisotropic-vectors} \hlt{Symmetries and Equal Quadratic Values}\\
Let $\mathcal{V}$ be nonsingular orthogonal with $\operatorname{char}(\mathbb{F})\neq2$, and put $q(\mathbf{x})\defeq\innerprod{\mathbf{x},\mathbf{x}}$.
\begin{enumerate}[label=(\roman*)]
\setlength{\itemsep}{0pt}
\item If $q(\mathbf{w})\neq0$, then $\sigma_{\mathbf{w}}$ is a linear involutive isometry, $\sigma_{\mathbf{w}}\mathbf{w}=-\mathbf{w}$, and $\sigma_{\mathbf{w}}|_{(\mathbb{F}\mathbf{w})^{\perp}}=I$.
\item If $q(\mathbf{u})=q(\mathbf{v})\neq0$, then at least one of $\mathbf{u}-\mathbf{v}$ and $\mathbf{u}+\mathbf{v}$ is nonisotropic; its symmetry maps $\mathbf{u}$ to $\mathbf{v}$ or $-\mathbf{v}$, respectively.
\end{enumerate}
\begin{proof}
\begin{enumerate}[label=(\roman*)]
\setlength{\itemsep}{0pt}
\item Fix $\mathbf{w}$ with $q(\mathbf{w})\neq0$. The formula is linear, sends $\mathbf{w}$ to $-\mathbf{w}$, and fixes $(\mathbb{F}\mathbf{w})^{\perp}$ pointwise.\\
By Theorem \ref{thm:metric-orthogonal-complement-dimension}(ii), $\mathcal{V}=\mathbb{F}\mathbf{w}\odot(\mathbb{F}\mathbf{w})^{\perp}$, so $\sigma_{\mathbf{w}}^2=I$ and $q\circ\sigma_{\mathbf{w}}=q$. Hence $\sigma_{\mathbf{w}}$ is bijective and isometric by Proposition \ref{prop:orthogonal-isometry-preserves-quadratic-form}.
\item Put $\mathbf{a}=\mathbf{u}-\mathbf{v}$ and $\mathbf{b}=\mathbf{u}+\mathbf{v}$. Then $\innerprod{\mathbf{a},\mathbf{b}}=q(\mathbf{u})-q(\mathbf{v})=0$.\\
If $q(\mathbf{a})=q(\mathbf{b})=0$, then $0=q(\mathbf{a}+\mathbf{b})=q(2\mathbf{u})=4q(\mathbf{u})$, a contradiction. Hence $q(\mathbf{a})\neq0$ or $q(\mathbf{b})\neq0$.\\
In the first case, $q(\mathbf{a})=2\innerprod{\mathbf{u},\mathbf{a}}$ and $\sigma_{\mathbf{a}}\mathbf{u}=\mathbf{u}-\mathbf{a}=\mathbf{v}$; in the second, $q(\mathbf{b})=2\innerprod{\mathbf{u},\mathbf{b}}$ and $\sigma_{\mathbf{b}}\mathbf{u}=\mathbf{u}-\mathbf{b}=-\mathbf{v}$.
\end{enumerate}
\hlt{Structural Commentary}: Verify one symmetry directly, then test the orthogonal directions $\mathbf{u}\pm\mathbf{v}$.\\
\hlt{Key Insight}: Equal nonzero quadratic values make at least one of the two difference directions reflective.
\end{proof}
\end{proposition}

\begin{theorem} \label{thm:orthogonal-transformations-products-of-symmetries} \hlt{Orthogonal Transformations Are Products of Symmetries}\\
Let $\mathcal{V}$ be nonsingular orthogonal with $\operatorname{char}(\mathbb{F})\neq2$, and put $q(\mathbf{x})\defeq\innerprod{\mathbf{x},\mathbf{x}}$. Then $\mathcal{O}(\mathcal{V})=\langle\sigma_{\mathbf{u}}:q(\mathbf{u})\neq0\rangle$.
\begin{proof}
Induct on $n=\dim\mathcal{V}$. Case $n=0$ is immediate. If $n=1$, every isometry acts by some $\lambda$ with $\lambda^2=1$, hence as $I$ or $\sigma_{\mathbf{u}}$.\\
Assume $n>1$. Nonsingularity provides $\mathbf{u}$ with $q(\mathbf{u})\neq0$; otherwise the form is alternating and zero by Proposition \ref{prop:form-type-relations-by-characteristic}(ii).\\
For $\tau\in\mathcal{O}(\mathcal{V})$, choose $\sigma$ by Proposition \ref{prop:symmetry-matches-equal-length-nonisotropic-vectors}(ii) with $\sigma\tau\mathbf{u}=\varepsilon\mathbf{u}$, where $\varepsilon\in\{1,-1\}$. Put $\rho=\sigma$ for $\varepsilon=1$ and $\rho=\sigma_{\mathbf{u}}\sigma$ for $\varepsilon=-1$; then $\rho\tau\mathbf{u}=\mathbf{u}$. If $\mathbf{x}\in(\mathbb{F}\mathbf{u})^{\perp}$, then $\innerprod{\rho\tau\mathbf{x},\mathbf{u}}=\innerprod{\rho\tau\mathbf{x},\rho\tau\mathbf{u}}=\innerprod{\mathbf{x},\mathbf{u}}=0$, so $\rho\tau$ preserves this complement.\\
It is nonsingular of dimension $n-1$ by Theorem \ref{thm:metric-orthogonal-complement-dimension}(i),(iv). The inductive factors of the restriction extend by $I$ on $\mathbb{F}\mathbf{u}$ to same-normal symmetries of $\mathcal{V}$. Thus $\rho\tau$ is a product of symmetries; so is $\rho^{-1}$ because each symmetry is involutive by Proposition \ref{prop:symmetry-matches-equal-length-nonisotropic-vectors}(i).\\
Hence $\tau=\rho^{-1}(\rho\tau)$ is such a product.\\
\hlt{Structural Commentary}: Fix one nonisotropic line by at most two symmetries, then induct on its orthogonal complement.\\
\hlt{Key Insight}: Symmetries play the role of elementary generators for the orthogonal group.
\end{proof}
\end{theorem}

\begin{proposition} \label{prop:isometry-fixed-space} \hlt{Fixed Spaces of Isometries}\\
Let $\mathcal{V}$ be nonsingular and let $\tau\in\operatorname{Iso}(\mathcal{V})$. Then $\ker(I-\tau)=\operatorname{im}(I-\tau)^{\perp}$.
\begin{proof}
If $\mathbf{x}\in\ker(I-\tau)$, then $(\forall\,\mathbf{y}\in\mathcal{V})\,\innerprod{\mathbf{x},(I-\tau)\mathbf{y}}=\innerprod{\mathbf{x},\mathbf{y}}-\innerprod{\tau\mathbf{x},\tau\mathbf{y}}=0$, so $\ker(I-\tau)\subseteq\operatorname{im}(I-\tau)^{\perp}$.\\
These spaces have equal dimension by Theorem \ref{thm:rank-nullity} and Theorem \ref{thm:metric-orthogonal-complement-dimension}(i), so they are equal.\\
\hlt{Structural Commentary}: Establish orthogonality of fixed vectors to displacements, then close the inclusion by dimension.\\
\hlt{Key Insight}: The fixed space is exactly the orthogonal complement of the displacement space.
\end{proof}
\end{proposition}

\begin{proposition} \label{prop:symmetries-fixed-hyperplanes} \hlt{Symmetries and Fixed Hyperplanes}\\
Let $\mathcal{V}$ be nonsingular orthogonal with $\operatorname{char}(\mathbb{F})\neq2$, put $q(\mathbf{x})\defeq\innerprod{\mathbf{x},\mathbf{x}}$, and let $\tau\in\mathcal{O}(\mathcal{V})$ with $\tau\neq I$.
\begin{enumerate}[label=(\roman*)]
\setlength{\itemsep}{0pt}
\item For nonisotropic $\mathbf{u}$, $\ker(I-\sigma_{\mathbf{u}})=(\mathbb{F}\mathbf{u})^{\perp}$ is a nonsingular hyperplane.
\item The condition $(\exists\,\mathbf{u}\in\mathcal{V})\,(q(\mathbf{u})\neq0\wedge\tau=\sigma_{\mathbf{u}})$ holds if and only if $\ker(I-\tau)$ is a nonsingular hyperplane.
\end{enumerate}
\begin{proof}
\begin{enumerate}[label=(\roman*)]
\setlength{\itemsep}{0pt}
\item[$(i)$] By Proposition \ref{prop:symmetry-matches-equal-length-nonisotropic-vectors}(i), $\sigma_{\mathbf{u}}$ fixes $(\mathbb{F}\mathbf{u})^{\perp}$ pointwise and sends $\mathbf{u}$ to $-\mathbf{u}$. Since $\operatorname{char}(\mathbb{F})\neq2$, its fixed space is exactly $(\mathbb{F}\mathbf{u})^{\perp}$, which is a nonsingular hyperplane by Theorem \ref{thm:metric-orthogonal-complement-dimension}(i),(iv).
\item[$(ii)$] ($\Rightarrow$) This is (i).
\item[$(ii)$] ($\Leftarrow$) Put $\mathcal{H}=\ker(I-\tau)$ and $\mathcal{L}=\mathcal{H}^{\perp}$. By Theorem \ref{hm:metric-orthogonal-complement-dimension}(ii),(iv), $\mathcal{V}=\mathcal{H}\odot\mathcal{L}$ with $\mathcal{L}$ a nonsingular line.\\
For $\mathbf{x}\in\mathcal{L}$ and $\mathbf{h}\in\mathcal{H}$, $\innerprod{\tau\mathbf{x},\mathbf{h}}=\innerprod{\tau\mathbf{x},\tau\mathbf{h}}=0$, so $\tau\mathcal{L}=\mathcal{L}$. Thus $\tau|_{\mathcal{L}}=\lambda I$ with $\lambda^2=1$.\\
Since $\tau\neq I$ and $\operatorname{char}(\mathbb{F})\neq2$, $\lambda=-1$. Every $\mathbf{0}\neq\mathbf{u}\in\mathcal{L}$ is nonisotropic, and $\tau=\sigma_{\mathbf{u}}$ by Proposition \ref{prop:symmetry-matches-equal-length-nonisotropic-vectors}(i).
\end{enumerate}
\hlt{Structural Commentary}: Read the fixed hyperplane from the symmetry formula, then recover the flipped normal line.\\
\hlt{Key Insight}: A nonidentity isometry is one symmetry exactly when its fixed space is a nonsingular hyperplane.
\end{proof}
\end{proposition}

\begin{proposition} \label{prop:anisotropic-space-has-no-nontrivial-unipotent-isometries} \hlt{No Unipotent Isometries in Anisotropic Spaces}\\
Let $\mathcal{V}$ be nonsingular anisotropic and let $\tau\in\operatorname{Iso}(\mathcal{V})$. If $(\exists\,m\in\mathbb{Z}_{>0})\,(\tau-I)^m=0$, then $\tau=I$.
\begin{proof}
Put $N=\tau-I$ and suppose $N\neq0$. Fix maximal $k\geq1$ with $N^k\neq0$, so $N^{k+1}=0$, and choose $\mathbf{x}$ with $\mathbf{u}\defeq N^k\mathbf{x}\neq\mathbf{0}$.\\
By Proposition \ref{prop:metric-adjoints-and-isometries}, $N^*=\tau^{-1}-I=-N+N^2-\cdots+(-1)^kN^k$.\\
Every monomial in $(N^*)^kN^k$ has $N$-degree at least $2k\geq k+1$, so this operator is zero.\\
Hence $q(\mathbf{u})=\innerprod{N^k\mathbf{x},N^k\mathbf{x}}=\innerprod{\mathbf{x},(N^*)^kN^k\mathbf{x}}=0$, contrary to anisotropy. Therefore $N=0$ and $\tau=I$.\\
\hlt{Structural Commentary}: Pass the top nilpotent displacement through the adjoint and force its quadratic value to vanish.\\
\hlt{Key Insight}: An anisotropic geometry has no nonzero isotropic displacement on which unipotence can terminate.
\end{proof}
\end{proposition}

Witt extension and cancellation apply here by Theorem \ref{thm:witt-extension-theorem-metric-spaces} and Corollary \ref{cor:witt-cancellation-metric-spaces}.
